% appendices/appendix_glossary.tex
\chapter*{附录A：术语表}
\addcontentsline{toc}{chapter}{附录A：术语表}

\small
\begin{longtable}{|p{0.24\textwidth}|p{0.41\textwidth}|p{0.25\textwidth}|}
\hline
\textbf{术语} & \textbf{定义} & \textbf{代码位置} \\
\hline
\endhead
BlockDevice & 表示任何可以读写固定大小块的存储介质的 trait 抽象 & src/block\_device.rs \\
dm-crypt & 提供透明磁盘加密的设备映射器目标 & src/dm\_crypt.rs \\
dm-verity & 提供透明完整性校验的设备映射器目标 & src/dm\_verity.rs \\
HashTree & dm-verity 使用的默克尔树结构 & src/hash\_tree.rs \\
AES-256-XTS & 磁盘加密的标准密码模式 & src/crypto.rs \\
plain64 IV & 使用 64 位扇区号作为 IV 的策略 & src/crypto.rs \\
Integrity\allowbreak Violation & 哈希不匹配时返回的错误 & src/error.rs \\
Root Hash & 默克尔树的顶层哈希，信任锚 & src/hash\_tree.rs \\
Tweak Key & XTS 密钥的后半部分，用于加密 IV & src/crypto.rs \\
IV Offset & 用于计算初始化向量的偏移值 & src/dm\_crypt.rs \\
Data Offset & 物理设备开头用于元数据的扇区数 & src/dm\_crypt.rs \\
DM\_SECTOR\_\allowbreak SIZE & 设备映射器的原子 I/O 单位（512 字节） & src/dm\_crypt.rs \\
DmDevice & 星绽内核中实现 BIO 映射的块设备 & aster-dm/device.rs \\
DM 启动参数 & crypt/verity 参数解析与 mapper 注册 & aster-dm/lib.rs \\
Passthrough\allowbreak Device & 直接转发调用的包装器 & device\_mapper.rs \\
aster-dm & 星绽内核中的 Device Mapper 组件 & kernel/comps/dm/ \\
DmTarget & 封装 Crypt 或 Verity 逻辑的枚举 & 补丁 0001 \\
self\_test\_\allowbreak ext2\_\allowbreak magic & 启动时自检函数 & 补丁 0002 \\
pivot\_root & 切换根文件系统的机制 & 补丁 0003 \\
dmctl & 用户态命令行工具 & dmctl/src/main.rs \\
MemoryBlock\allowbreak Device & 基于内存的后端 & src/block\_device.rs \\
FileBlock\allowbreak Device & 基于文件的后端 & src/block\_device.rs \\
\hline
\end{longtable}
